\section{引言：微观何时可以被遗忘}
\label{sec:introduction}

复杂系统的宏观描述通常做两件事：遗忘大量微观差异，并用较少的状态预测后续变化。前者是压缩，后者是自治。二者缺一不可。若只要求状态更少，常值映射已经给出最强压缩，却通常不保留任务；若只要求转移闭合，恒等表示永远精确闭合，却没有简化。本文研究的问题是：

\begin{quote}
在任务、微观探针、允许动力学和一步时间尺度固定以后，什么时候能够遗忘绝大多数操作上真实可区分的微观信息，同时保留任务，并让宏观状态近似自行演化？
\end{quote}

这是一项条件判定问题，不是一条“所有复杂系统必有简单规律”的存在宣言。本文采用操作主义的最低限度立场：可表达、可检验的规律依赖于对经验情形的区分、标记和概率关系；数学模型还必须明确哪些区分属于当前任务与系统边界。这个立场只解释为何需要任务和探针，并不替代后文的闭合条件。

\subsection{问题的两个陷阱}

第一个陷阱是\emph{编码填充}。给每个状态附加一个不影响任务和动力学的复制标签，就能任意放大原始状态数。因此，$|Y|/|X|$ 很小本身没有表示不变的意义。本文先由预先声明的微观探针和允许动力学构造操作性微观空间 $B$，剔除严格不可区分的复制；宏观表示必须再是 $B$ 的任务保真严格信息商。状态数只在这个信息顺序确定后才用于衡量压缩。

第二个陷阱是\emph{把候选表示的定义当成存在证明}。如果先假设有一个低复杂度、近似闭合的映射，再推出它对应一条宏观规律，结论主要是重述假设。本文的中心结果改为一个不含候选宏观转移核的充要判据：对任务保真分区 $\pi$，同时检查分区的相对大小和同一块内状态的下一块预测分歧；二者能否同时趋于零，恰好刻画系统族是否存在“少而自治”的表示。

\subsection{与既有工作的关系}

给定分区后，有限 Markov 链何时能够精确聚合，是 lumpability 理论的经典问题\citep{KemenySnell1976,Buchholz1994}；Markov functions 与线性不变性给出了相关的投影条件\citep{RogersPitman1981,GurvitsLedoux2005}。概率互模拟与状态空间最小化从行为等价和算法细化的角度研究相近结构\citep{LarsenSkou1991,DerisaviEtAl2003}。计算力学中的因果状态则把未来预测等价作为状态构造原则\citep{CrutchfieldYoung1989,ShaliziCrutchfield2001}。本文不把这些已有结论重新命名为“涌现定理”；新增的组织重点是：

\begin{enumerate}
  \item 在比较宏观与微观状态数以前构造操作性微观基线，并证明其对严格无意义复制填充不变；
  \item 以纤维预测直径定义系统族联合指标 $\Gamma_N$，证明 $\Gamma_N\to0$ 与任务相对低复杂度近似闭合等价；
  \item 给出贯穿全文的非平凡系统族：可交换故障--修复链的宏观状态数从 $2^N$ 精确降为且最少需要 $N+1$；
  \item 区分候选的闭合判定与候选构造，并用残余对称和 Transformer 预测状态说明后者。
\end{enumerate}

Markov category 为随机映射及其复合提供了统一语言\citep{Fritz2020}，但范畴化并非本文主定理的逻辑前提；相关组织放在附录~\ref{app:categories}。标准 Borel 空间、随机粗粒化和连续时间也需要额外的可测性、因子化或生成元条件，本文只在第~\ref{sec:discussion}节与附录~\ref{app:extensions}中给出边界清楚的推广。

\subsection{结论强度与路线}

本文的“自治”专指给定动力学族下的一步 Markov 闭合；“简单”专指相对于操作性微观空间的状态类别较少。结论不自动蕴含公式短、参数少、计算便宜、语义可解释、因果稳定、低维流形、跨分布泛化或长时间误差小。全文依赖链只有四步：

\[
\text{操作性微观基线}
\longrightarrow
\text{纤维半径--直径引理}
\longrightarrow
\Gamma_N\text{ 充要判据}
\longrightarrow
\text{具体模型证书}.
\]

第~\ref{sec:operational}节定义信息层级；第~\ref{sec:closure}节建立闭合的可检验桥梁；第~\ref{sec:criterion}节给出中心定理；第~\ref{sec:fault-repair}节证明非平凡正例；第~\ref{sec:construction}节讨论候选表示从何而来；第~\ref{sec:verification}节集中说明可计算性、估计稳健性和结论边界。独立证明与扩展全部移入附录和补充材料。
